\chapter{2013年湖南卷（文）}

\begin{problem}
在平面直角坐标系 \(xOy\) 中，若直线 \(l_{1}: \left\{ \begin{array}{l} x = 2s + 1, \\ y = s, \end{array} \right.\) （\(s\) 为参数）和直线 \(l_{2}: \left\{ \begin{array}{l} x = at, \\ y = 2t - 1, \end{array} \right.\) （\(t\) 为参数）平行，则常数 \(a\) 的值为\fillinblank{}.
\end{problem}

\section{单选题}

\begin{problem}
复数 \(z = \i \cdot (1 + \i)\) (i 为虚数单位) 在复平面上对应的点位于（\quad）

\choices
    {第一象限}
    {第二象限}
    {第三象限}
    {第四象限}
\end{problem}

\begin{problem}
“1 < x < 2”是“x < 2”成立的（\quad）

\choices
    {充分不必要条件}
    {必要不充分条件}
    {充分必要条件}
    {既不充分也不必要条件}
\end{problem}

\begin{problem}
某工厂甲，乙，丙三个车间生产了同一种产品，数量分别为 120 件，80 件，60 件. 为了解它们的产品质量是否存在显著差异，用分层抽样方法抽取了一个容量为 n 的样本进行调查，其中从丙车间的产品中抽取了 3 件，则 n=（\quad）

\choices
    {9}
    {10}
    {12}
    {13}
\end{problem}

\begin{problem}
已知 \( f(x) \) 是奇函数，\( g(x) \) 是偶函数，且 \( f(-1)+g(1)=2 \) ，\( f(1)+g(-1)=4 \) ，则 \( g(1) \) 等于（\quad）

\choices
    {4}
    {3}
    {2}
    {1}
\end{problem}

\begin{problem}
在锐角 \(\triangle ABC\) 中，角 \(A, B\) 所对的边长分别为 \(a, b\). 若 \(2a \sin B = \sqrt{3} b\)，则角 \(A\) 等于（\quad）

\choices
    {\( \frac{\pi}{3} \)}
    {\( \frac{\pi}{4} \)}
    {\( \frac{\pi}{6} \)}
    {\( \frac{\pi}{12} \)}
\end{problem}

\begin{problem}
函数 \( f(x) = \ln x \) 的图象与函数 \( g(x) = x^2 - 4x + 4 \) 的图象的交点个数为（\quad）

\choices
    {0}
    {1}
    {2}
    {3}
\end{problem}

\begin{problem}
已知正方体的棱长为 1，其俯视图是一个面积为 1 的正方形，侧视图是一个面积为 \(\sqrt{2}\) 的矩形，则该正方体的正视图的面积等于（\quad）

\choices
    {\(\frac{\sqrt{3}}{2}\)}
    {1}
    {\(\frac{\sqrt{2} + 1}{2}\)}
    {\(\sqrt{2}\)}
\end{problem}

\begin{problem}
已知 \(a, b\) 是单位向量，\(a \cdot b = 0\). 若向量 \(c\) 满足 \(|c - a - b| = 1\)，则 \(|c|\) 的最大值为（\quad）

\choices
    {\(\sqrt{2} - 1\)}
    {\(\sqrt{2}\)}
    {\(\sqrt{2} + 1\)}
    {\(\sqrt{2} + 2\)}
\end{problem}

\begin{problem}
已知事件“在矩形 \(ABCD\) 的边 \(CD\) 上随机取一点 \(P\)，使 \(\triangle APB\) 的最大边是 \(AB\)”发生的概率为 \(\frac{1}{2}\)，则 \(\frac{AD}{AB} =\)（\quad）

\choices
    {\(\frac{1}{2}\)}
    {\(\frac{1}{4}\)}
    {\(\frac{\sqrt{3}}{2}\)}
    {\(\frac{\sqrt{7}}{4}\)}
\end{problem}

\section{填空题}

\begin{problem}
已知集合 \(U = \{2, 3, 6, 8\}\)，\(A = \{2, 3\}\)，\(B = \{2, 6, 8\}\)，则 \(\complement_U A \cap B =\fillinblank{}\).
\end{problem}

\begin{problem}
在平面直角坐标系 \(xOy\) 中，若直线 \(l_{1}: \left\{ \begin{array}{l} x = 2s + 1, \\ y = s, \end{array} \right.\) （\(s\) 为参数）和直线 \(l_{2}: \left\{ \begin{array}{l} x = at, \\ y = 2t - 1, \end{array} \right.\) （\(t\) 为参数）平行，则常数 \(a\) 的值为\fillinblank{}.
\end{problem}

\begin{problem}
执行如图所示的程序框图，如果输入 \(a = 1, b = 2\) ，则输出的 \(a\) 的值为\fillinblank{}.

\bitmapfigure[max width=0.46\linewidth,max totalheight=4.0cm]{img/2013/hunan_liberal/q12_fig1.png}
\end{problem}

\begin{problem}
若变量 \(x, y\) 满足约束条件 \(\left\{ \begin{array}{l} x + 2y \leqslant 8, \\ 0 \leqslant x \leqslant 4, \\ 0 \leqslant y \leqslant 3, \end{array} \right.\) 则 \(x + y\) 的最大值为\fillinblank{}.
\end{problem}

\begin{problem}
设 \(F_{1}, F_{2}\) 是双曲线 \(C: \frac{x^{2}}{a^{2}} - \frac{y^{2}}{b^{2}} = 1 (a > 0, b > 0)\) 的两个焦点. 若在 \(C\) 上存在一点 \(P\)，使 \(PF_{1} \perp PF_{2}\)，且 \(\angle PF_{1}F_{2} = 30^{\circ}\)，则 \(C\) 的离心率为\fillinblank{}.
\end{problem}

\begin{problem}
对于 \(E = \{a_1, a_2, \cdots, a_{100}\}\) 的子集 \(X = \{a_{i1}, a_{i2}, \cdots, a_{ik}\}\)，定义 \(X\) 的“特征数列”为 \(x_1, x_2, \cdots, x_{100}\)，其中 \(x_{i1} = x_{i2} = \cdots = x_{ik} = 1\)，其余项均为 0，例如：子集 \(\{a_2, a_3\}\) 的“特征数列”为 \(0, 1, 1, 0, 0, \cdots, 0\).
\begin{enumerate}
\item 子集 \( \{a_{1}, a_{3}, a_{5}\} \) 的“特征数列”的前 3 项和等于\fillinblank{}.
\item 若 E 的子集 \(P\) 的“特征数列” \(p_{1}, p_{2}, \cdots, p_{100}\) 满足 \(p_{1} = 1, p_{i} + p_{i+1} = 1, 1 \leqslant i \leqslant 99\)；E 的子集 \(Q\) 的“特征数列” \(q_{1}, q_{2}, \cdots, q_{100}\) 满足 \(q_{1} = 1, q_{j} + q_{j+1} + q_{j+2} = 1, 1 \leqslant j \leqslant 98\)，则 \(P \cap Q\) 的元素个数为\fillinblank{}.
\end{enumerate}
\end{problem}

\section{解答题}

\begin{problem}
已知函数 \( f(x) = \cos x \cdot \cos \left(x - \frac{\pi}{3}\right) \).
\begin{enumerate}
\item 求 \( f\left(\frac{2\pi}{3}\right) \) 的值；
\item 求使 \( f(x) < \frac{1}{4} \) 成立的 \( x \) 的取值集合.
\end{enumerate}
\end{problem}

\begin{problem}
如图，在直棱柱 \(ABC - A_{1}B_{1}C_{1}\) 中，\(\angle BAC = 90^{\circ}\)，\(AB = AC = \sqrt{2}\)，\(AA_{1} = 3, D\) 是 \(BC\) 的中点，点 \(E\) 在棱 \(BB_{1}\) 上运动.

\begin{enumerate}
\item 证明： \( AD \perp C_{1}E; \)
\item 当异面直线 \(AC, C_1E\) 所成的角为 \(60^\circ\) 时，求三棱锥 \(C_1 - A_1B_1E\) 的体积.
\end{enumerate}

\bitmapfigure[max width=0.42\linewidth,max totalheight=4.0cm]{img/2013/hunan_liberal/q17_fig1.png}
\end{problem}

\begin{problem}
某人在如图所示的直角边长为 4 米的三角形地块的每个格点 (指纵，横直线的交叉点以及三角形的顶点) 处都种了一株相同品种的作物. 根据历年的种植经验，一株该种作物的年收获量 \(Y\) (单位: \( \mathrm{kg} \)) 与它的“相近”作物株数 \(X\) 之间的关系如下表所示：X 1 2 3 4 Y 51 48 45 42 这里，两株作物“相近”是指它们之间的直线距离不超过 1 米.

\begin{enumerate}
\item 完成下表，并求所种作物的平均年收获量； Y 51 48 45 42 频数 4
\item 在所种作物中随机选取一株，求它的年收获量至少为 \(48 \mathrm{~kg}\) 的概率.
\end{enumerate}

\bitmapfigure[max width=0.36\linewidth,max totalheight=4.0cm]{img/2013/hunan_liberal/q18_fig1.png}
\end{problem}

\begin{problem}
设 \(S_{n}\) 为数列 \(\{a_{n}\}\) 的前 \(n\) 项和. 已知 \(a_1 \neq 0, 2a_{n} - a_1 = S_1 \cdot S_n, n \in \mathbf{N}^*\).
\begin{enumerate}
\item 求 \( a_{1}, a_{2} \) ，并求数列 \( \{a_{n}\} \) 的通项公式；
\item 求数列 \( \{na_{n}\} \) 的前 n 项和.
\end{enumerate}
\end{problem}

\begin{problem}
已知 \(F_{1}, F_{2}\) 分别是椭圆 \(E: \frac{x^{2}}{5} + y^{2} = 1\) 的左，右焦点，\(F_{1}, F_{2}\) 关于直线 \(x + y - 2 = 0\) 的对称点是圆 \(C\) 的一条直径的两个端点.
\begin{enumerate}
\item 求圆 C 的方程；
\item 设过点 \(F_{2}\) 的直线 \(l\) 被椭圆 \(E\) 和圆 \(C\) 所截得的弦长分别为 \(a, b\). 当 \(ab\) 最大时，求直线 \(l\) 的方程.
\end{enumerate}
\end{problem}

\begin{problem}
已知函数 \( f(x) = \frac{1 - x}{1 + x^{2}} \e^x \).
\begin{enumerate}
\item 求 \( f(x) \) 的单调区间；
\item 证明：当 \( f(x_{1}) = f(x_{2}) (x_{1} \neq x_{2}) \) 时，\( x_{1} + x_{2} < 0 \) .
\end{enumerate}
\end{problem}
